\documentclass[12pt]{article}

% Idioma y codificación
\usepackage[spanish, es-tabla, es-notilde]{babel}       %es-tabla para que se titule "Tabla"
\usepackage[utf8]{inputenc}

% Márgenes
\usepackage[a4paper,top=3cm,bottom=2.5cm,left=3cm,right=3cm]{geometry}

% Comentarios de bloque
\usepackage{verbatim}

% Paquetes de links
\usepackage[hidelinks]{hyperref}    % Permite enlaces
\usepackage{url}                    % redirecciona a la web

% Más opciones para enumeraciones
\usepackage{enumitem}

% Personalizar la portada
\usepackage{titling}

% Paquetes de tablas
\usepackage{multirow}

% Para añadir el símbolo de euro
\usepackage{eurosym}

% Paquete para subrayar con diferentes estilos
\usepackage[normalem]{ulem}


%------------------------------------------------------------------------

%Paquetes de figuras
\usepackage{caption}
\usepackage{subcaption} % Figuras al lado de otras
\usepackage{float}      % Poner figuras en el sitio indicado H.


% Paquetes de imágenes
\usepackage{graphicx}       % Paquete para añadir imágenes
\usepackage{transparent}    % Para manejar la opacidad de las figuras

% Paquete para usar colores
\usepackage[dvipsnames, table, xcdraw]{xcolor}
\usepackage{pagecolor}      % Para cambiar el color de la página

% Habilita tamaños de fuente mayores
\usepackage{fix-cm}

% Para los gráficos
\usepackage{tikz}
\usepackage{forest}

% Para poder situar los nodos en los grafos
\usetikzlibrary{positioning}

% Para evitar problemas entre babel y tikz
\usetikzlibrary{babel}


%------------------------------------------------------------------------

% Paquetes de matemáticas
\usepackage{mathtools, amsfonts, amssymb, mathrsfs}
\usepackage[makeroom]{cancel}     % Simplificar tachando
\usepackage{polynom}    % Divisiones y Ruffini
\usepackage{units} % Para poner fracciones diagonales con \nicefrac

\usepackage{pgfplots}   %Representar funciones
\pgfplotsset{compat=1.18}  % Versión 1.18

\usepackage{tikz-cd}    % Para usar diagramas de composiciones
\usetikzlibrary{calc}   % Para usar cálculo de coordenadas en tikz

%Definición de teoremas, etc.
\usepackage{amsthm}
%\swapnumbers   % Intercambia la posición del texto y de la numeración

\theoremstyle{plain}

\makeatletter
\@ifclassloaded{article}{
  % --- Configuración para ARTICLE ---
  \newtheorem{teo}{Teorema}[section]
  \theoremstyle{definition}
  \newtheorem{definicion}{Definición}[section]
  \newtheorem{ejercicio}{Ejercicio}[section]

  \renewcommand{\theteo}{%
    \ifnum\value{section}=0 \arabic{teo}\else \thesection.\arabic{teo}\fi
  }
  \renewcommand{\thedefinicion}{%
    \ifnum\value{section}=0 \arabic{definicion}\else \thesection.\arabic{definicion}\fi
  }
  \renewcommand{\theejercicio}{%
    \ifnum\value{section}=0 \arabic{ejercicio}\else \thesection.\arabic{ejercicio}\fi
  }
}{
  % --- Configuración para BOOK / REPORT ---
  \newtheorem{teo}{Teorema}[chapter]
  \theoremstyle{definition}
  \newtheorem{definicion}{Definición}[chapter]
  \newtheorem{ejercicio}{Ejercicio}[section]

  % Si no hay sección: Capítulo.Teorema (ej. 2.1)
  % Si hay sección:    Capítulo.Sección.Teorema (ej. 2.3.1)
  \renewcommand{\theteo}{%
    \ifnum\value{section}=0
      \thechapter.\arabic{teo}%
    \else
      \thesection.\arabic{teo}% (\thesection ya contiene el número de capítulo)
    \fi
  }
  \renewcommand{\thedefinicion}{%
    \ifnum\value{section}=0
      \thechapter.\arabic{definicion}%
    \else
      \thesection.\arabic{definicion}%
    \fi
  }
  \renewcommand{\theejercicio}{%
    \ifnum\value{section}=0
      \arabic{ejercicio}%
    \else
      \thesection.\arabic{ejercicio}% Formato por defecto: ej. 2.2.1
    \fi
  }
}
\makeatother

% Teoremas dependientes o con contador compartido
\newtheorem{coro}{Corolario}[teo]
\newtheorem{prop}[teo]{Proposición}
\newtheorem{lema}[teo]{Lema}
\newtheorem*{lema*}{Lema}

\theoremstyle{remark}
\newtheorem*{observacion}{Observación}

\theoremstyle{definition}
\newtheorem*{notacion}{Notación}
\newtheorem*{ejemplo}{Ejemplo}
\newtheorem*{ejercicio*}{Ejercicio}


% \renewcommand\qedsymbol{$\blacksquare$}         % Cambiar símbolo QED
%------------------------------------------------------------------------

% Paquetes para encabezados
\usepackage{fancyhdr}
\pagestyle{fancy}
\fancyhf{}

\newcommand{\helv}{ % Modificación tamaño de letra
\fontfamily{}\fontsize{12}{12}\selectfont}
\setlength{\headheight}{15pt} % Amplía el tamaño del índice


%\usepackage{lastpage}   % Referenciar última pag   \pageref{LastPage}
%\fancyfoot[C]{%
%  \begin{minipage}{\textwidth}
%    \centering
%    ~\\
%    \thepage\\
%    \href{https://losdeldgiim.github.io/}{\texttt{\footnotesize losdeldgiim.github.io}}
%  \end{minipage}
%}
\fancyfoot[C]{\thepage}
\fancyfoot[R]{\href{https://losdeldgiim.github.io/}{\texttt{\footnotesize losdeldgiim.github.io}}}

%------------------------------------------------------------------------

% Conseguir que no ponga "Capítulo 1". Sino solo "1."
\makeatletter
\@ifclassloaded{book}{
  \renewcommand{\chaptermark}[1]{\markboth{\thechapter.\ #1}{}} % En el encabezado
    
  \renewcommand{\@makechapterhead}[1]{%
  \vspace*{50\p@}%
  {\parindent \z@ \raggedright \normalfont
    \ifnum \c@secnumdepth >\m@ne
      \huge\bfseries \thechapter.\hspace{1em}\ignorespaces
    \fi
    \interlinepenalty\@M
    \Huge \bfseries #1\par\nobreak
    \vskip 40\p@
  }}
}
\makeatother

%------------------------------------------------------------------------
% Paquetes de cógido
\usepackage{minted}
\renewcommand\listingscaption{Código fuente}

\usepackage{fancyvrb}
% Personaliza el tamaño de los números de línea
\renewcommand{\theFancyVerbLine}{\small\arabic{FancyVerbLine}}

% Estilo para C++
\newminted{cpp}{
    frame=lines,
    framesep=2mm,
    baselinestretch=1.2,
    linenos,
    escapeinside=||
}

% para minted
\definecolor{LightGray}{rgb}{0.95,0.95,0.92}
\setminted{
    linenos=true,
    stepnumber=5,
    numberfirstline=true,
    autogobble,
    breaklines=true,
    breakautoindent=true,
    breaksymbolleft=,
    breaksymbolright=,
    breaksymbolindentleft=0pt,
    breaksymbolindentright=0pt,
    breaksymbolsepleft=0pt,
    breaksymbolsepright=0pt,
    fontsize=\footnotesize,
    bgcolor=LightGray,
    numbersep=10pt
}


\usepackage{listings} % Para incluir código desde un archivo

\renewcommand\lstlistingname{Código Fuente}
\renewcommand\lstlistlistingname{Índice de Códigos Fuente}

% Definir colores
\definecolor{vscodepurple}{rgb}{0.5,0,0.5}
\definecolor{vscodeblue}{rgb}{0,0,0.8}
\definecolor{vscodegreen}{rgb}{0,0.5,0}
\definecolor{vscodegray}{rgb}{0.5,0.5,0.5}
\definecolor{vscodebackground}{rgb}{0.97,0.97,0.97}
\definecolor{vscodelightgray}{rgb}{0.9,0.9,0.9}

% Configuración para el estilo de C similar a VSCode
\lstdefinestyle{vscode_C}{
  backgroundcolor=\color{vscodebackground},
  commentstyle=\color{vscodegreen},
  keywordstyle=\color{vscodeblue},
  numberstyle=\tiny\color{vscodegray},
  stringstyle=\color{vscodepurple},
  basicstyle=\scriptsize\ttfamily,
  breakatwhitespace=false,
  breaklines=true,
  captionpos=b,
  keepspaces=true,
  numbers=left,
  numbersep=5pt,
  showspaces=false,
  showstringspaces=false,
  showtabs=false,
  tabsize=2,
  frame=tb,
  framerule=0pt,
  aboveskip=10pt,
  belowskip=10pt,
  xleftmargin=10pt,
  xrightmargin=10pt,
  framexleftmargin=10pt,
  framexrightmargin=10pt,
  framesep=0pt,
  rulecolor=\color{vscodelightgray},
  backgroundcolor=\color{vscodebackground},
}

%------------------------------------------------------------------------
% La etiqueda IA sirve para indicar a partir de donde le pasamos el preambulo a la IA para el escaneo
% IA

% Comandos definidos
\newcommand{\bb}[1]{\mathbb{#1}}
\newcommand{\cc}[1]{\mathcal{#1}}

% I prefer the slanted \leq
\let\oldleq\leq % save them in case they're every wanted
\let\oldgeq\geq
\renewcommand{\leq}{\leqslant}
\renewcommand{\geq}{\geqslant}

% Si y solo si
\newcommand{\sii}{\iff}

% MCD y MCM
\DeclareMathOperator{\mcd}{mcd}
\DeclareMathOperator{\mcm}{mcm}

% Signo
\DeclareMathOperator{\sgn}{sgn}

% Imagen
\DeclareMathOperator{\Ima}{Im}

% Aplicación identidad
\newcommand{\id}{\mathrm{Id}}

% Letras griegas
\newcommand{\eps}{\epsilon}
\newcommand{\veps}{\varepsilon}
\newcommand{\lm}{\lambda}

\newcommand{\ol}{\overline}
\newcommand{\ul}{\underline}
\newcommand{\wt}{\widetilde}
\newcommand{\wh}{\widehat}

\let\oldvec\vec
\renewcommand{\vec}{\overrightarrow}

% Derivadas parciales
\newcommand{\del}[2]{\frac{\partial #1}{\partial #2}}
\newcommand{\Del}[3]{\frac{\partial^{#1} #2}{\partial #3^{#1}}}
\newcommand{\deld}[2]{\dfrac{\partial #1}{\partial #2}}
\newcommand{\Deld}[3]{\dfrac{\partial^{#1} #2}{\partial #3^{#1}}}


\newcommand{\AstIg}{\stackrel{(\ast)}{=}}
\newcommand{\Hop}{\stackrel{L'H\hat{o}pital}{=}}

\newcommand{\red}[1]{{\color{red}#1}} % Para integrales, destacar los cambios.

% Método de integración
\newcommand{\MetInt}[2]{
    \left[\begin{array}{c}
        #1 \\ #2
    \end{array}\right]
}

% Declarar aplicaciones
% 1. Nombre aplicación
% 2. Dominio
% 3. Codominio
% 4. Variable
% 5. Imagen de la variable
\newcommand{\Func}[5]{
    \begin{equation*}
        \begin{array}{rrll}
            \displaystyle #1:& \displaystyle  #2 & \longrightarrow & \displaystyle  #3\\
               & \displaystyle  #4 & \longmapsto & \displaystyle  #5
        \end{array}
    \end{equation*}
}

%------------------------------------------------------------------------


\newcommand{\myhref}[2]{\href{#1}{\color{blue}\uline{#2}}}
\newcommand{\T}[0]{\cc{T}}

\begin{document}

    % 1. Foto de fondo
    % 2. Título
    % 3. Encabezado Izquierdo
    % 4. Color de fondo
    % 5. Coord x del titulo
    % 6. Coord y del titulo
    % 7. Fecha

    
    % 1. Foto de fondo
% 2. Título
% 3. Encabezado Izquierdo
% 4. Color de fondo
% 5. Coord x del titulo
% 6. Coord y del titulo
% 7. Fecha
% 8. Autor

\newcommand{\portada}[8]{

    \portadaBase{#1}{#2}{#3}{#4}{#5}{#6}{#7}{#8}
    \portadaBook{#1}{#2}{#3}{#4}{#5}{#6}{#7}{#8}
}

\newcommand{\portadaExamen}[8]{

    \portadaBase{#1}{#2}{#3}{#4}{#5}{#6}{#7}{#8}
    \portadaArticle{#1}{#2}{#3}{#4}{#5}{#6}{#7}{#8}
}




\newcommand{\portadaBase}[8]{

    % Tiene la portada principal y la licencia Creative Commons
    
    % 1. Foto de fondo
    % 2. Título
    % 3. Encabezado Izquierdo
    % 4. Color de fondo
    % 5. Coord x del titulo
    % 6. Coord y del titulo
    % 7. Fecha
    % 8. Autor
    
    
    \thispagestyle{empty}               % Sin encabezado ni pie de página
    \newgeometry{margin=0cm}        % Márgenes nulos para la primera página
    
    
    % Encabezado
    \fancyhead[L]{\helv #3}
    \fancyhead[R]{\helv \nouppercase{\leftmark}}
    
    
    \pagecolor{#4}        % Color de fondo para la portada
    
    \begin{figure}[p]
        \centering
        \transparent{0.3}           % Opacidad del 30% para la imagen
        
        \includegraphics[width=\paperwidth, keepaspectratio]{../../_assets/#1}
    
        \begin{tikzpicture}[remember picture, overlay]
            \node[anchor=north west, text=white, opacity=1, font=\fontsize{60}{90}\selectfont\bfseries\sffamily, align=left] at (#5, #6) {#2};
            
            \node[anchor=south east, text=white, opacity=1, font=\fontsize{12}{18}\selectfont\sffamily, align=right] at (9.7, 3) {\href{https://losdeldgiim.github.io/}{\textbf{Los Del DGIIM}, \texttt{\footnotesize losdeldgiim.github.io}}};
            
            \node[anchor=south east, text=white, opacity=1, font=\fontsize{12}{15}\selectfont\sffamily, align=right] at (9.7, 1.8) {Doble Grado en Ingeniería Informática y Matemáticas\\Universidad de Granada};
        \end{tikzpicture}
    \end{figure}
    
    
    \restoregeometry        % Restaurar márgenes normales para las páginas subsiguientes
    \pagecolor{white}       % Restaurar el color de página
    
    
    \newpage
    \thispagestyle{empty}               % Sin encabezado ni pie de página
    \begin{tikzpicture}[remember picture, overlay]
        \node[anchor=south west, inner sep=3cm] at (current page.south west) {
            \begin{minipage}{0.5\paperwidth}
                \href{https://creativecommons.org/licenses/by-nc-nd/4.0/}{
                    \includegraphics[height=2cm]{../../_assets/Licencia.png}
                }\vspace{1cm}\\
                Esta obra está bajo una
                \href{https://creativecommons.org/licenses/by-nc-nd/4.0/}{
                    Licencia Creative Commons Atribución-NoComercial-SinDerivadas 4.0 Internacional (CC BY-NC-ND 4.0).
                }\\
    
                Eres libre de compartir y redistribuir el contenido de esta obra en cualquier medio o formato, siempre y cuando des el crédito adecuado a los autores originales y no persigas fines comerciales. 
            \end{minipage}
        };
    \end{tikzpicture}
    
    
    
    % 1. Foto de fondo
    % 2. Título
    % 3. Encabezado Izquierdo
    % 4. Color de fondo
    % 5. Coord x del titulo
    % 6. Coord y del titulo
    % 7. Fecha
    % 8. Autor


}


\newcommand{\portadaBook}[8]{

    % 1. Foto de fondo
    % 2. Título
    % 3. Encabezado Izquierdo
    % 4. Color de fondo
    % 5. Coord x del titulo
    % 6. Coord y del titulo
    % 7. Fecha
    % 8. Autor

    % Personaliza el formato del título
    \pretitle{\begin{center}\bfseries\fontsize{42}{56}\selectfont}
    \posttitle{\par\end{center}\vspace{2em}}
    
    % Personaliza el formato del autor
    \preauthor{\begin{center}\Large}
    \postauthor{\par\end{center}\vfill}
    
    % Personaliza el formato de la fecha
    \predate{\begin{center}\huge}
    \postdate{\par\end{center}\vspace{2em}}
    
    \title{#2}
    \author{\href{https://losdeldgiim.github.io/}{Los Del DGIIM, \texttt{\large losdeldgiim.github.io}}
    \\ \vspace{0.5cm}#8}
    \date{Granada, #7}
    \maketitle
    
    \tableofcontents
}




\newcommand{\portadaArticle}[8]{

    % 1. Foto de fondo
    % 2. Título
    % 3. Encabezado Izquierdo
    % 4. Color de fondo
    % 5. Coord x del titulo
    % 6. Coord y del titulo
    % 7. Fecha
    % 8. Autor

    % Personaliza el formato del título
    \pretitle{\begin{center}\bfseries\fontsize{42}{56}\selectfont}
    \posttitle{\par\end{center}\vspace{2em}}
    
    % Personaliza el formato del autor
    \preauthor{\begin{center}\Large}
    \postauthor{\par\end{center}\vspace{3em}}
    
    % Personaliza el formato de la fecha
    \predate{\begin{center}\huge}
    \postdate{\par\end{center}\vspace{5em}}
    
    \title{#2}
    \author{\href{https://losdeldgiim.github.io/}{Los Del DGIIM, \texttt{\large losdeldgiim.github.io}}
    \\ \vspace{0.5cm}#8}
    \date{Granada, #7}
    \thispagestyle{empty}               % Sin encabezado ni pie de página
    \maketitle
    \vfill
}
    \portadaExamen{ffccA4.jpg}{Topología I\\Examen XVII}{Topología I. Examen XVII}{MidnightBlue}{-8}{28}{2026}{Lucía Benjumeda Maldonado}

    \begin{description}
        \item[Asignatura] Topología I.
        \item[Curso Académico] 2026-27.
        \item[Grado] Grado en Matemáticas.
        \item[Grupo] A.
        \item[Profesor] Ana María Hurtado Cortegana.
        \item[Descripción] Parcial.
        \item[Fecha] 18 de noviembre de 2026.
        % \item[Duración] 60 minutos.
    
    \end{description}
    \newpage

    \begin{ejercicio} Defina el concepto de \emph{base} de una topología $\cc{T}$ sobre un conjunto $X$.
    Sea $(X, \cc{T})$ un espacio topológico, $\cc{B}$ una base de $\cc{T}$ y $A \subseteq X$ un subconjunto no vacío. Demuestre que la familia
    \[
        \cc{B}_A = \{ B \cap A : B \in \cc{B} \}
    \]
    constituye una base para la topología inducida $\cc{T}_A$ sobre $A$.
    \end{ejercicio}

    \begin{ejercicio}
    En el plano real $\bb{R}^2$, considérese la familia de subconjuntos
    \[
        \cc{B} = \big\{ \{x\} \times [y, +\infty) : x, y \in \bb{R} \big\}.
    \]
    Se pide:
    \begin{enumerate}[label=\alph*)]
        \item Demuestre que $\cc{B}$ es base de una topología $T$ en $\bb{R}^2$.
        
        \item Determine qué rectas del plano $\bb{R}^2$ heredan la topología discreta como subespacios topológicos de $(\bb{R}^2, T)$.
        
        \item Calcule el interior, la clausura y la frontera en $(\bb{R}^2, T)$ de los siguientes subconjuntos:
        \begin{enumerate}[label=\roman*)]
            \item $D = \{ (x, y) \in \bb{R}^2 : x^2 + y^2 < 1 \}$.
            \item $Q = \{0\} \times \bb{Q}$.
            \item $H = [0, +\infty) \times \bb{R}$.
        \end{enumerate}
        
        \item Estudie si el espacio topológico $(\bb{R}^2, T)$ verifica las propiedades de separación $T_1$ y $T_2$ (Hausdorff), así como el primer axioma de numerabilidad ($1\mathrm{AN}$) y el segundo axioma de numerabilidad ($2\mathrm{AN}$).
    \end{enumerate}
    \end{ejercicio}

    \newpage
    \setcounter{ejercicio}{0} % Reiniciar contador de ejercicios
    \noindent
\begin{ejercicio}
Defina el concepto de \emph{base} de una topología $\cc{T}$ sobre un conjunto $X$.
Sea $(X, \cc{T})$ un espacio topológico, $\cc{B}$ una base de $\cc{T}$ y $A \subseteq X$ un subconjunto no vacío. Demuestre que la familia
\[
    \cc{B}_A = \{ B \cap A : B \in \cc{B} \}
\]
constituye una base para la topología inducida $\cc{T}_A$ sobre $A$.

\medskip
Llamamos \emph{base} de una topología $\cc{T}$ sobre un conjunto $X$ a una familia de abiertos $\cc{B} \subseteq \cc{T}$ (llamados \emph{abiertos básicos}) tal que todo abierto $U \in \cc{T}$ puede expresarse como unión de elementos de $\cc{B}$; es decir, existe una subfamilia $\{B_i\}_{i \in I} \subseteq \cc{B}$ tal que
\[
    U = \bigcup_{i \in I} B_i,
\]
o, equivalentemente, para todo $U \in \cc{T}$ y todo $x \in U$, existe $B \in \cc{B}$ tal que $x \in B \subseteq U$.

Sea ahora $(X, \cc{T})$ un espacio topológico, $\cc{B}$ una base de $\cc{T}$ y $A \subseteq X$ un subconjunto no vacío. Recordemos que la topología inducida sobre $A$ viene dada por $\cc{T}_A = \{ U \cap A : U \in \cc{T} \}$. Para probar que $\cc{B}_A$ es base de $\cc{T}_A$, comprobamos las dos condiciones de la definición:
\begin{itemize}
    \item $\cc{B}_A \subseteq \cc{T}_A$: Dado $B \cap A \in \cc{B}_A$, como $B \in \cc{B} \subseteq \cc{T}$, por definición de topología inducida se tiene directamente que $B \cap A \in \cc{T}_A$.
    \item Todo abierto de $\cc{T}_A$ es unión de elementos de $\cc{B}_A$: Sea $V \in \cc{T}_A$ arbitrario. Existe $U \in \cc{T}$ tal que $V = U \cap A$. Como $\cc{B}$ es base de $\cc{T}$, podemos escribir $\displaystyle U = \bigcup_{i \in I} B_i$ con $B_i \in \cc{B}$ para cada $i \in I$. Aplicando la distributividad de la intersección respecto de la unión:
    \[
        V = U \cap A = \left( \bigcup_{i \in I} B_i \right) \cap A = \bigcup_{i \in I} (B_i \cap A).
    \]
    Como cada término $B_i \cap A$ pertenece a $\cc{B}_A$, queda demostrado que $\cc{B}_A$ es base de $\cc{T}_A$.
\end{itemize}
\end{ejercicio}

\vspace{0.5cm}

\begin{ejercicio}
En el plano real $\bb{R}^2$, considérese la familia de subconjuntos
\[
    \cc{B} = \big\{ \{x\} \times [y, +\infty) : x, y \in \bb{R} \big\}.
\]
Se pide:
\begin{enumerate}[label=\alph*)]
    \item Demuestre que $\cc{B}$ es base de una topología $\cc{T}$ en $\bb{R}^2$.
    
    Veamos que la familia $\cc{B}$ cumple las dos condiciones caracterizadoras de base:
    \begin{itemize}
        \item Por doble inclusión veremos que $\displaystyle \bigcup_{B \in \cc{B}} B = \bb{R}^2$:
        \begin{itemize}
            \item[($\subseteq$)] Dado $(x, y) \in \displaystyle \bigcup_{B \in \cc{B}} B$, existe $B = \{x\} \times [y_0, +\infty) \in \cc{B}$ con $y_0 \leq y$ tal que $(x, y) \in B \subseteq \bb{R}^2$.
            \item[($\supseteq$)] Dado $(x, y) \in \bb{R}^2$, basta tomar $B = \{x\} \times [y, +\infty) \in \cc{B}$. Se cumple $(x, y) \in B$, luego $(x, y) \in \displaystyle \bigcup_{B \in \cc{B}} B$.
        \end{itemize}
        \item Sean $B_1 = \{x_1\} \times [y_1, +\infty)$ y $B_2 = \{x_2\} \times [y_2, +\infty)$ dos elementos de $\cc{B}$ y supongamos que $(x, y) \in B_1 \cap B_2$. Entonces $x = x_1 = x_2$ e $y \geq \max(y_1, y_2)$. Tomando $y_3 = \max(y_1, y_2)$, el conjunto $B_3 = \{x\} \times [y_3, +\infty)$ pertenece a $\cc{B}$ y cumple
        \[
            (x, y) \in B_3 \subseteq B_1 \cap B_2.
        \]
    \end{itemize}
    Por el teorema de caracterización de bases, existe una única topología $\cc{T}(\cc{B})$ en $\bb{R}^2$ que tiene a $\cc{B}$ como base.

    \item Determine qué rectas del plano $\bb{R}^2$ heredan la topología discreta como subespacios topológicos de $(\bb{R}^2, \cc{T}(\cc{B}))$.

    Una recta $L \subseteq \bb{R}^2$ hereda la topología discreta $\cc{T}_{|L}$ si y solo si todo conjunto unitario $\{(x_0, y_0)\} \subseteq L$ es abierto en $\cc{T}_{|L}$; es decir, si para cada $(x_0, y_0) \in L$ existe un abierto básico $B \in \cc{B}$ tal que $B \cap L = \{(x_0, y_0)\}$. Distinguimos dos casos:
    \begin{itemize}
        \item Rectas no verticales: ($L : y = mx + n$, con $m, n \in \bb{R}$): Dado $(x_0, y_0) \in L$, elegimos el básico $B = \{x_0\} \times [y_0, +\infty) \in \cc{B}$. Como la recta $L$ corta a cada vertical $\{x_0\} \times \bb{R}$ en un único punto, se tiene
        \[
            B \cap L = \big(\{x_0\} \times [y_0, +\infty)\big) \cap L = \{(x_0, y_0)\}.
        \]
        Por tanto, todos los puntos son abiertos en $\cc{T}_{|L}$ y las rectas no verticales sí son discretas.
        \item Rectas verticales ($L = \{x_0\} \times \bb{R}$, con $x_0 \in \bb{R}$): Para cualquier $(x_0, y_0) \in L$, todo abierto $U \in \cc{T}$ que contenga a $(x_0, y_0)$ contiene algún básico $B = \{x_0\} \times [y', +\infty)$ con $y' \leq y_0$. Entonces
        \[
            B \cap L = \{x_0\} \times [y', +\infty) = B \neq \{(x_0, y_0)\},
        \]
        ya que contiene infinitos puntos de $L$ por encima de $y_0$. Así, ningún punto es abierto en $L$ y las rectas verticales no son discretas.
    \end{itemize}

    \begin{center}
    \begin{tikzpicture}[scale=0.85]
        % Ejes
        \draw[->, gray] (-2.2,0) -- (2.5,0) node[right] {$x$};
        \draw[->, gray] (0,-1.5) -- (0,2.5) node[above] {$y$};
        
        % Recta no vertical
        \draw[blue, thick] (-2,-1.2) -- (2.2,1.6) node[right] {$L$};
        % Semirrecta básica
        \draw[red, very thick, ->] (0.8, 0.67) -- (0.8, 2.3) node[above] {$B$};
        \fill[red] (0.8, 0.67) circle (2.2pt) node[below right, black] {\small $(x_0,y_0)$};
        \node[align=center, font=\footnotesize] at (0.2, -2.0) {Subespacio discreto};

        % Segundo panel: recta vertical
        \begin{scope}[xshift=6.5cm]
            \draw[->, gray] (-1.5,0) -- (2.5,0) node[right] {$x$};
            \draw[->, gray] (0,-1.5) -- (0,2.5) node[above] {$y$};
            % Recta vertical
            \draw[blue, thick] (1.0,-1.4) -- (1.0,2.4) node[right] {$L$};
            % Semirrecta básica
            \draw[red, very thick, dashed, ->] (1.0, 0.2) -- (1.0, 2.3) node[above left] {$B$};
            \fill[red] (1.0, 0.5) circle (2.2pt) node[left, black] {\small $(x_0,y_0)$};
            \node[align=center, font=\footnotesize] at (0.6, -2.0) {No discreto};
        \end{scope}
    \end{tikzpicture}
    \end{center}

    \item Calcule el interior, la clausura y la frontera en $(\bb{R}^2, \cc{T})$ de los siguientes subconjuntos:
    \begin{enumerate}
        \item $D = \{ (x, y) \in \bb{R}^2 : x^2 + y^2 < 1 \}$.\\
        $D^{\circ} = \varnothing$.\\
        Dado cualquier $(x, y) \in D$, todo abierto básico que lo contenga es de la forma $B = \{x\} \times [y', +\infty)$ con $y' \leq y$. Como $B$ no es acotado superiormente en la segunda coordenada y $D$ sí es acotado, se cumple $B \not\subseteq D$. Luego $(x, y) \notin D^{\circ}$ y $D^{\circ} = \varnothing$.\\

        $\overline{D} = D \cup \big((-1, 1) \times (-\infty, 0]\big)$.\\
        Probamos la igualdad por doble inclusión:
        \begin{itemize}
            \item[($\supseteq$)] Como $D \subseteq \overline{D}$, basta ver que todo $(x, y) \in C \setminus D$ pertenece a $\overline{D}$. Si $(x, y) \in C \setminus D$, entonces $|x| < 1$ e $y \leq -\sqrt{1 - x^2} < 0$. Dado cualquier básico $B = \{x\} \times [y', +\infty)$ que contenga a $(x, y)$, se cumple $y' \leq y < 0 < \sqrt{1 - x^2}$, por lo que $(x, 0) \in B \cap D \neq \varnothing$. Así, $(x, y) \in \overline{D}$.
            \item[($\subseteq$)] Sea $(x, y) \in \bb{R}^2 \setminus C$. Distinguimos dos casos:
            \begin{enumerate}
                \item Si $|x| \geq 1$, el básico $B = \{x\} \times [y, +\infty)$ contiene a $(x, y)$ y es disjunto con $D$ (pues todo punto de $D$ tiene $|x| < 1$).
                \item Si $|x| < 1$ e $y \geq \sqrt{1 - x^2}$, el básico $B = \{x\} \times [y, +\infty)$ contiene a $(x, y)$ y cumple $B \cap D = \varnothing$ (pues todos sus puntos satisfacen $x^2 + \tilde{y}^2 \geq x^2 + y^2 \geq 1$).
            \end{enumerate}
            En ambos casos $(x, y) \notin \overline{D}$, de donde $\overline{D} \subseteq C$.
        \end{itemize}
        $\partial D = \overline{D} \setminus D^{\circ} = \overline{D} = D \cup \big((-1, 1) \times (-\infty, 0]\big)$.
        \begin{center}
        \begin{tikzpicture}[scale=1]
            % Panel 1: Clausura de D
            \begin{scope}
                \fill[blue!20] (-1,-2.2) rectangle (1,0);
                \fill[blue!20] (-1,0) arc (180:0:1) -- cycle;
                \draw[dashed, blue, thick] (-1,-2.2) -- (-1,0) arc (180:0:1) -- (1,-2.2);
                \draw[dotted, gray] (-1,0) arc (180:360:1);
                \draw[->, gray] (-1.8,0) -- (1.8,0) node[right] {$x$};
                \draw[->, gray] (0,-2.3) -- (0,1.6) node[above] {$y$};
                \node[font=\footnotesize] at (0, -2.7) {$\overline{D}$};
            \end{scope}
        \end{tikzpicture}
        \end{center}
        
        \item $Q = \{0\} \times \bb{Q}$.\\
        $Q^{\circ} = \varnothing$.\\
        Sea $(0, q) \in Q$. Cualquier abierto que lo contenga incluye un básico $B = \{0\} \times [y', +\infty)$ con $y' \leq q$. Por la densidad de los irracionales en $\bb{R}$, existe $p \in (\bb{R} \setminus \bb{Q}) \cap [y', +\infty)$, de modo que $(0, p) \in B$ pero $(0, p) \notin Q$. Así, $B \not\subseteq Q$ y $(0, q) \notin Q^{\circ}$.\\

        $\overline{Q} = \{0\} \times \bb{R}$.\\
        Si $x \neq 0$, el básico $B = \{x\} \times [y, +\infty)$ contiene a $(x, y)$ y verifica $B \cap Q = \varnothing$, luego $(x, y) \notin \overline{Q}$. Por el contrario, si $x = 0$, para cualquier $y \in \bb{R}$ y cualquier básico $B = \{0\} \times [y', +\infty)$ con $y' \leq y$, por densidad de $\bb{Q}$ en $\bb{R}$ existe $q \in \bb{Q} \cap [y', +\infty)$, por lo que $(0, q) \in B \cap Q \neq \varnothing$. Por tanto, $\overline{Q} = \{0\} \times \bb{R}$.\\

        $\partial Q = \overline{Q} \setminus Q^{\circ} = \{0\} \times \bb{R}$.
        \begin{center}
        \begin{tikzpicture}
        % Panel 2: Q y clausura de Q
        \begin{scope}[xshift=5.2cm]
            \draw[->, gray] (-1.8,0) -- (1.8,0) node[right] {$x$};
            \draw[->, gray] (0,-2.3) -- (0,1.6) node[above] {$y$};
            \draw[red, very thick] (0,-2.2) -- (0,1.4);
            \foreach \y in {-2.0,-1.6,...,1.2} {
                \fill[blue] (0,\y) circle (1.5pt);
            }
            \node[align=center, font=\footnotesize] at (0, -2.7) {$Q$};
        \end{scope}
        \end{tikzpicture}
        \end{center}

        \item $H = [0, +\infty) \times \bb{R}$.\\
        $H^{\circ} = H$.\\
        Para cada $(x, y) \in H$ se tiene $x \geq 0$, luego el básico $B = \{x\} \times [y, +\infty)$ satisface $(x, y) \in B \subseteq H$. Así, $H$ es abierto y $H^{\circ} = H$.\\

        $\overline{H} = H$.\\
        Su complementario es $\bb{R}^2 \setminus H = (-\infty, 0) \times \bb{R}$. Para todo $(x, y) \in \bb{R}^2 \setminus H$ se cumple $x < 0$, luego $B = \{x\} \times [y, +\infty) \subseteq (-\infty, 0) \times \bb{R}$. Por tanto, $\bb{R}^2 \setminus H$ es abierto, $H$ es cerrado y $\overline{H} = H$.\\

        $\partial H = \overline{H} \setminus H^{\circ} = \varnothing$.
        \begin{center}
        \begin{tikzpicture}
        \begin{scope}[xshift=10.4cm]
                \fill[green!20] (0,-2.2) rectangle (1.8,1.5);
                \draw[green!60!black, very thick] (0,-2.2) -- (0,1.5);
                \draw[->, gray] (-1.5,0) -- (2.0,0) node[right] {$x$};
                \draw[->, gray] (0,-2.3) -- (0,1.6) node[above] {$y$};
                \node[align=center, font=\footnotesize] at (0.2, -2.7) {$H$};
            \end{scope}
        \end{tikzpicture}
        \end{center}
    \end{enumerate}

    

    \item Estudie si el espacio topológico $(\bb{R}^2, \cc{T})$ verifica las propiedades de separación $T_1$ y $T_2$ (Hausdorff), así como el primer axioma de numerabilidad ($1\mathrm{AN}$) y el segundo axioma de numerabilidad ($2\mathrm{AN}$).
    \begin{itemize}
        \item \textbf{Axiomas de separación ($T_1$ y $T_2$):} No es $T_1$ ni, por consiguiente, $T_2$.\\
        Tomemos dos puntos distintos en la misma vertical $(x, y), (x, y') \in \bb{R}^2$ con $y < y'$. Todo abierto $U \in \cc{T}$ tal que $(x, y) \in U$ contiene al básico $\{x\} \times [y, +\infty)$. Como $y < y'$, se tiene $(x, y') \in \{x\} \times [y, +\infty) \subseteq U$. Por tanto, no existe ningún abierto que contenga a $(x, y)$ y deje fuera a $(x, y')$, lo que impide que $(\bb{R}^2, \cc{T})$ sea $T_1$ (y tampoco $T_2$).

        \begin{center}
        \begin{tikzpicture}
            \draw[->, gray] (-1.5,0) -- (2.2,0) node[right] {$x$};
            \draw[->, gray] (0,-0.8) -- (0,2.5) node[above] {$y$};
            \draw[red, very thick, ->] (0.9, 0.4) -- (0.9, 2.3) node[right] {$B_{(x,y)}$};
            \fill[black] (0.9, 0.4) circle (2.2pt) node[left] {\small $(x,y)$};
            \fill[blue] (0.9, 1.5) circle (2.2pt) node[left] {\small $(x,y')$};
            \node[align=center, font=\footnotesize] at (0.4, -1.3) {};
        \end{tikzpicture}
        \end{center}

        \item \textbf{Primer axioma de numerabilidad ($1\mathrm{AN}$):} Sí lo verifica.\\
        Para cada punto $(x, y) \in \bb{R}^2$, la familia unitaria $\cc{B}_{(x, y)} = \big\{ \{x\} \times [y, +\infty) \big\}$ es una base de entornos finita (luego numerable) de $(x, y)$, ya que dado cualquier abierto $U \in \cc{T}$ con $(x, y) \in U$, existe un básico $\{x\} \times [y', +\infty) \subseteq U$ con $y' \leq y$, y por tanto $\{x\} \times [y, +\infty) \subseteq \{x\} \times [y', +\infty) \subseteq U$.

        \item \textbf{Segundo axioma de numerabilidad ($2\mathrm{AN}$):} No lo verifica.\\
        Por el apartado b), la recta horizontal $L = \bb{R} \times \{0\}$ es un subespacio discreto no numerable (tiene el cardinal de $\bb{R}$). Si $(\bb{R}^2, \cc{T})$ fuese $2\mathrm{AN}$, esta propiedad sería hereditaria y el subespacio $(L, \cc{T}_L)$ debería admitir una base numerable. Sin embargo, un espacio discreto solo es $2\mathrm{AN}$ si es numerable, lo cual es una contradicción.
    \end{itemize}
\end{enumerate}
\end{ejercicio}
\end{document}